\begin{table}[htbp]
\centering
\small
\setlength{\tabcolsep}{4pt}
\caption{C-seq batch by batch: the intended count per bin and the deviation, achieved minus intended, then the cumulative deviation after batches 2 and 3 (cum.); mean ± sd over r1-r5, r0 (exploratory) apart. Each batch reallocates against what the earlier batches actually returned; the cumulative columns show whether retraining corrects a deviation or moves it. The outside rows count points whose true FOS left the selection range, where nothing was intended (generator bins, outer bin FOS > 6). None of this says what a perfect implementation would do; that is the oracle's question. This table is r0 alone: exploratory, never averaged in. C0-log, C0-seq and C0-logG train on the same guided points as C-log, C-seq and C-logG in every realization, so their allocation, batch and displacement rows are those of C, shared rather than duplicated.}
\label{tab:structural:batches:r0}
\begin{tabular}{lrrrrrrrr}
\toprule
bin & int. 1 & dev. 1 & int. 2 & dev. 2 & int. 3 & dev. 3 & cum. 2 & cum. 3 \\
\midrule
0.50-0.59 & 22 & -7 & 13 & -6 & 12 & -2 & -13 & -15 \\
0.59-0.70 & 19 & -1 & 7 & 2 & 4 & -1 & 1 & 0 \\
0.70-0.82 & 16 & -6 & 12 & 0 & 6 & -4 & -6 & -10 \\
0.82-0.97 & 13 & -8 & 14 & -5 & 11 & -1 & -13 & -14 \\
0.97-1.14 & 7 & -2 & 7 & 2 & 5 & 1 & 0 & 1 \\
1.14-1.35 & 3 & 0 & 5 & -1 & 7 & 2 & -1 & 1 \\
1.35-1.59 & 0 & 2 & 0 & 3 & 0 & 1 & 5 & 6 \\
1.59-1.88 & 0 & 0 & 2 & 0 & 6 & -1 & 0 & -1 \\
1.88-2.22 & 1 & -1 & 6 & -2 & 8 & -1 & -3 & -4 \\
2.22-2.62 & 0 & 0 & 0 & 1 & 4 & 2 & 1 & 3 \\
2.62-3.09 & 0 & 1 & 0 & 0 & 6 & 0 & 1 & 1 \\
3.09-3.65 & 0 & 1 & 1 & 1 & 5 & 3 & 2 & 5 \\
3.65-4.31 & 2 & -2 & 8 & -3 & 9 & -5 & -5 & -10 \\
4.31-5.08 & 1 & 0 & 6 & -2 & 8 & -2 & -2 & -4 \\
5.08-6.00 & 16 & -13 & 19 & -3 & 9 & 1 & -16 & -15 \\
< 0.5 (outside) & 0 & 22 & 0 & 6 & 0 & 4 & 28 & 32 \\
> 6 (outside) & 0 & 14 & 0 & 7 & 0 & 3 & 21 & 24 \\
\bottomrule
\end{tabular}
\end{table}
